% Chapter of the map: Information and Kolmogorov complexity.
% Conventions and the template of an entry: blueprint/README.md. Check with
% blueprint/tools/check_names.sh. Labels of this file begin with kind:information (def:information-..., thm:information-...).
\chapter{Information and Kolmogorov complexity}
\label{chap:information}

This chapter collects the results of the map about description length, the charge for running
time, and the conservation laws that relate them. It starts with two measures of the information in
a word, the shortest program that prints it and the same with a charge for the running time, in
two models (the partial recursive codes of Mathlib and the machines of cslib), and with the law
that a computation does not create description length. It continues with logical depth, with
counting laws for a solver, a hint and a verifier in a model of query trees over finite sets, with
a closed universe in which the verifier and the justification of a ``no'' are paid for, and with
conservation of work in a black-box model. It ends with Blum speed-up, the cut of costs, and two
small facts about indices and determinants. In the models of codes the cost is the step budget of
\texttt{evaln}, which is not the running time of a machine; in the black-box models there is no
machine, and the statements say which model they are in.

\section{Kolmogorov complexity and its time-charged version}
\label{sec:information-kolmogorov}

In this section a program is a partial recursive code of Mathlib. The length $|p|$ of a code is the
length of the list of bits that spells out its tree in preorder, three bits for each node. A code
$p$ \emph{produces} $x$ when $p$ on input $0$ halts with value $x$. The cost $\mathrm{cost}(p, n)$ is
the least budget $k$ for which \texttt{evaln} halts on $p$ and $n$; it is a measure of Blum, not a
number of machine steps. Also, \texttt{correrMT}$(i, w, t)$ is the output, if any, of the cslib
machine numbered $i$ on the word $w$ after $t$ steps of the universal function of cslib.

\begin{definition}[Kolmogorov complexity and Levin's $\mathrm{Kt}$, on codes of Mathlib]
  \label{def:information-k-kt}
  \lean{Sagitta.Map.Tiempo.KolmogorovTiempo.K, Sagitta.Map.Tiempo.KolmogorovTiempo.Kt}
  \leanok
  For $x \in \mathbb{N}$, $\KC(x)$ is the least length $|p|$ of a code $p$ that produces $x$, and
  $\mathrm{Kt}(x)$ is the least value of $|p| + \lfloor \log_2 \mathrm{cost}(p, 0) \rfloor$ over the
  codes $p$ that produce $x$.
  \source{\cite{Kolmogorov65, Levin73}.}
  \fidelity{Not checked against the printed source.}
  \leanfile{Sagitta/Map/KolmogorovTiempo.lean}
\end{definition}

\begin{theorem}[$\KC \le \mathrm{Kt}$, on codes of Mathlib]
  \label{thm:information-k-le-kt}
  \lean{Sagitta.Map.Tiempo.KolmogorovTiempo.K_le_Kt}
  \leanok
  \uses{def:information-k-kt}
  For every $x \in \mathbb{N}$, $\KC(x) \le \mathrm{Kt}(x)$.
  \source{\cite{Kolmogorov65, Levin73}.}
  \fidelity{Not checked against the printed source.}
  \leanfile{Sagitta/Map/KolmogorovTiempo.lean}
\end{theorem}

\begin{proof}
  \leanok
  Both measures are minima of $|p| + g(\mathrm{cost}(p,0))$ over the programs that produce $x$, with
  $g = 0$ and $g = \lfloor \log_2 \cdot \rfloor$; a minimizing program for the second one gives a
  value at least $|p| \ge \KC(x)$.
\end{proof}

\begin{definition}[Kolmogorov complexity and $\mathrm{Kt}$, on machines of cslib]
  \label{def:information-machine-k}
  \lean{Sagitta.Map.Informacion.NoCreacion.ImprimeMT, Sagitta.Map.Informacion.NoCreacion.Kmt, Sagitta.Map.Informacion.NoCreacion.Ktmt}
  \leanok
  % uses-lean: Sagitta.Map.Tiempo.LevinMaquinas.correrMT
  The machine numbered $i$ \emph{prints} the word $x \in \{0,1\}^*$ with clock $t$ when
  $\texttt{correrMT}(i, [\,], t) = \mathrm{some}\ x$. Then $\mathrm{Kmt}(x)$ is the least $i$ for
  which some $t$ makes machine $i$ print $x$ with clock $t$, and $\mathrm{Ktmt}(x)$ is the least
  value of $i + \lceil \log_2 t \rceil$ over the pairs $(i, t)$ such that machine $i$ prints $x$
  with clock $t$.
  \source{\cite{Kolmogorov65, Levin73}.}
  \fidelity{Not checked against the printed source.}
  \leanfile{Sagitta/Map/Informacion/NoCreacion.lean}
\end{definition}

\begin{theorem}[The minimal machine and the minimal pair print the word]
  \label{thm:information-kmt-spec}
  \lean{Sagitta.Map.Informacion.NoCreacion.Kmt_spec, Sagitta.Map.Informacion.NoCreacion.Ktmt_spec}
  \leanok
  \uses{def:information-machine-k}
  For every word $x \in \{0,1\}^*$ there is $t$ such that the machine numbered $\mathrm{Kmt}(x)$
  prints $x$ with clock $t$. For every word $x$ there are $i$ and $t$ such that machine $i$ prints
  $x$ with clock $t$ and $i + \lceil \log_2 t \rceil = \mathrm{Ktmt}(x)$.
  \source{\cite{Kolmogorov65, Levin73}.}
  \fidelity{Not checked against the printed source.}
  \leanfile{Sagitta/Map/Informacion/NoCreacion.lean}
\end{theorem}

\begin{proof}
  \leanok
  Some machine prints every word, namely the constant machine of cslib, so the sets over which the
  two measures take a minimum are not empty; the minimum of a nonempty set of natural numbers is
  attained.
\end{proof}

\section{No creation of information}
\label{sec:information-no-creation}

A computation does not create description length: the program for $f(x)$ is the program for $x$
followed by the program for $f$. The statements about codes are in the model of Mathlib codes of
the previous section; the last one is about cslib machines and assumes a hypothesis on composition.

\begin{theorem}[No creation of information, on codes of Mathlib]
  \label{thm:information-no-creation}
  \lean{Sagitta.Map.Informacion.NoCreacion.no_creacion}
  \leanok
  \uses{def:information-k-kt}
  Let $p$ be a code, $f : \mathbb{N} \to \mathbb{N}$ and $x \in \mathbb{N}$, and assume that $p$
  on input $x$ halts with value $f(x)$. Then $\KC(f(x)) \le \KC(x) + |p| + 3$.
  \own{}
  \leanfile{Sagitta/Map/Informacion/NoCreacion.lean}
\end{theorem}

\begin{proof}
  \leanok
  Compose $p$ with a shortest program for $x$. The composition produces $f(x)$, and its length is
  the sum of the two lengths plus the three bits of the composition node.
\end{proof}

\begin{theorem}[No creation of information, with the time charge, on codes of Mathlib]
  \label{thm:information-no-creation-kt}
  \lean{Sagitta.Map.Informacion.NoCreacion.no_creacion_Kt}
  \leanok
  \uses{def:information-k-kt}
  Let $p$ be a code, $f : \mathbb{N} \to \mathbb{N}$ and $x \in \mathbb{N}$, and assume that $p$
  on input $x$ halts with value $f(x)$. Then
  \[
    \mathrm{Kt}(f(x)) \le \mathrm{Kt}(x) + \bigl(|p| + \lfloor \log_2 \mathrm{cost}(p, x) \rfloor\bigr) + 3 .
  \]
  \own{}
  \leanfile{Sagitta/Map/Informacion/NoCreacion.lean}
\end{theorem}

\begin{proof}
  \leanok
  Compose $p$ with a program that realizes $\mathrm{Kt}(x)$. The \texttt{evaln} budget of the
  composition is at most the maximum of the two budgets, and the logarithm of a maximum is at most
  the sum of the logarithms.
\end{proof}

\begin{theorem}[No creation of information, with no time charge, for a total function]
  \label{thm:information-no-crece}
  \lean{Sagitta.Map.Tiempo.ProfundidadLogica.K_no_crece}
  \leanok
  \uses{def:information-k-kt}
  Let $q$ be a code and $f : \mathbb{N} \to \mathbb{N}$ such that $q$ computes $f$, that is,
  $q$ on input $n$ halts with value $f(n)$ for every $n$. Then for every $x \in \mathbb{N}$,
  $\KC(f(x)) \le \KC(x) + |q| + 3$, where $3$ is the constant \texttt{costoComponer}.
  \own{}
  \leanfile{Sagitta/Map/ProfundidadLogica.lean}
\end{theorem}

\begin{proof}
  \leanok
  The statement for one input applies the shortest program for $x$ composed with $q$; the
  constant $3$ is the length of the label of the composition node.
\end{proof}

\begin{definition}[Hypothesis: a composer of machine indices]
  \label{def:information-compositor}
  \lean{Sagitta.Map.Informacion.NoCreacion.Compositor}
  \leanok
  For $g : \mathbb{N} \to \mathbb{N} \to \mathbb{N}$ and $c \in \mathbb{N}$, the proposition
  $\mathrm{Compositor}(g, c)$ says: for all $i, j, t, s$ and words $x, y$, if
  $\texttt{correrMT}(i, [\,], t) = \mathrm{some}\ x$ and $\texttt{correrMT}(j, x, s) =
  \mathrm{some}\ y$, then $\texttt{correrMT}(g(i,j), [\,], t + s + c) = \mathrm{some}\ y$. The
  entry only says that the definition exists; no entry asserts it.
  \own{}
  \leanfile{Sagitta/Map/Informacion/NoCreacion.lean}
\end{definition}

\begin{theorem}[No creation of information, on cslib machines, given a composer]
  \label{thm:information-no-creation-mt}
  \lean{Sagitta.Map.Informacion.NoCreacion.no_creacion_mt}
  \leanok
  \uses{def:information-compositor, def:information-machine-k}
  Let $g : \mathbb{N} \to \mathbb{N} \to \mathbb{N}$ and $c \in \mathbb{N}$, and assume
  $\mathrm{Compositor}(g, c)$. Let $f$ be a function on words, $j, s \in \mathbb{N}$ and $x$ a word,
  and assume $\texttt{correrMT}(j, x, s) = \mathrm{some}\ f(x)$. Then
  $\mathrm{Kmt}(f(x)) \le g(\mathrm{Kmt}(x), j)$.
  \own{}
  \leanfile{Sagitta/Map/Informacion/NoCreacion.lean}
\end{theorem}

\begin{proof}
  \leanok
  Machine $\mathrm{Kmt}(x)$ prints $x$ with some clock; the composer gives a machine $g(\mathrm{Kmt}(x), j)$
  that prints $f(x)$, and $\mathrm{Kmt}$ is at most the index of any machine that prints.
\end{proof}

\section{Logical depth and computational depth}
\label{sec:information-depth}

Both depths below are in the model of Mathlib codes, with $\KC$, $\mathrm{Kt}$, length and cost
as in Section~\ref{sec:information-kolmogorov}. Reading the input of a code already costs
\texttt{evaln} budget, so none of these results is a statement about machine time.

\begin{definition}[Computational depth, on codes of Mathlib]
  \label{def:information-prof}
  \lean{Sagitta.Map.Informacion.Profundidad.prof}
  \leanok
  \uses{def:information-k-kt}
  For $x \in \mathbb{N}$, $\mathrm{prof}(x) = \mathrm{Kt}(x) - \KC(x)$, with the subtraction of
  natural numbers.
  \source{\cite{Bennett88, Antunes06}.}
  \fidelity{Not checked against the printed source.}
  \leanfile{Sagitta/Map/Informacion/Profundidad.lean}
\end{definition}

\begin{definition}[Logical depth with significance $s$, on codes of Mathlib]
  \label{def:information-logical-depth}
  \lean{Sagitta.Map.Tiempo.ProfundidadLogica.profundidadLogica}
  \leanok
  \uses{def:information-k-kt}
  For $s, x \in \mathbb{N}$, the logical depth $\mathrm{depth}_s(x)$ is the least value of
  $\mathrm{cost}(p, 0)$ over the codes $p$ that produce $x$ and satisfy $|p| \le \KC(x) + s$.
  \source{\cite{Bennett88}.}
  \fidelity{Not checked against the printed source.}
  \leanfile{Sagitta/Map/ProfundidadLogica.lean}
\end{definition}

\begin{theorem}[Logical depth does not increase with the significance, on codes of Mathlib]
  \label{thm:information-depth-antitone}
  \lean{Sagitta.Map.Tiempo.ProfundidadLogica.profundidadLogica_antitona}
  \leanok
  \uses{def:information-logical-depth}
  For every $x \in \mathbb{N}$, the function $s \mapsto \mathrm{depth}_s(x)$ is antitone: if
  $s \le s'$ then $\mathrm{depth}_{s'}(x) \le \mathrm{depth}_s(x)$.
  \source{\cite{Bennett88}.}
  \fidelity{Not checked against the printed source.}
  \leanfile{Sagitta/Map/ProfundidadLogica.lean}
\end{theorem}

\begin{proof}
  \leanok
  A program admitted for the significance $s$ is also admitted for $s'$, so the minimum over the
  larger set is not larger.
\end{proof}

\begin{theorem}[Slow growth of logical depth, weak form, on codes of Mathlib]
  \label{thm:information-slow-growth-weak}
  \lean{Sagitta.Map.Tiempo.ProfundidadLogica.crecimiento_lento}
  \leanok
  \uses{def:information-logical-depth}
  Let $q$ be a code, $x, y \in \mathbb{N}$ and $s \in \mathbb{N}$, and assume that $q$ on input $x$
  halts with value $y$. Then
  \[
    \mathrm{depth}_{\KC(x) + s + |q| + 3 - \KC(y)}(y) \le \max\bigl(\mathrm{depth}_s(x),\ \mathrm{cost}(q, x)\bigr),
  \]
  with the subtraction of natural numbers.
  \source{\cite{Bennett88}.}
  \fidelity{Not checked against the printed source.}
  \leanfile{Sagitta/Map/ProfundidadLogica.lean}
\end{theorem}

\begin{proof}
  \leanok
  Compose $q$ with a program that attains $\mathrm{depth}_s(x)$. By no creation the composition is
  admitted at the significance of the statement, and the \texttt{evaln} budget of a composition is
  at most the maximum of the two budgets.
\end{proof}

\begin{theorem}[Slow growth of computational depth, with a computable inverse, on codes of Mathlib]
  \label{thm:information-slow-growth-inverse}
  \lean{Sagitta.Map.Informacion.Profundidad.crecimiento_lento}
  \leanok
  \uses{def:information-prof, thm:information-no-creation, thm:information-no-creation-kt}
  Let $p, r$ be codes, $f : \mathbb{N} \to \mathbb{N}$ and $x \in \mathbb{N}$, and assume that $p$
  on input $x$ halts with value $f(x)$ and $r$ on input $f(x)$ halts with value $x$. Then
  \[
    \mathrm{prof}(f(x)) \le \mathrm{prof}(x) + \lfloor \log_2 \mathrm{cost}(p, x) \rfloor + (|p| + |r| + 6).
  \]
  \source{\cite{Antunes06}.}
  \fidelity{Not checked against the printed source.}
  \leanfile{Sagitta/Map/Informacion/Profundidad.lean}
\end{theorem}

\begin{proof}
  \leanok
  Add the bound of no creation with the time charge for $\mathrm{Kt}(f(x))$ and the bound of no
  creation for $\KC(x)$ in terms of $\KC(f(x))$ through $r$, and use $\mathrm{prof} + \KC =
  \mathrm{Kt}$.
\end{proof}

\begin{definition}[Hypothesis: a code takes no longer than a finite deterministic system]
  \label{def:information-simula}
  \lean{Sagitta.Map.Informacion.MemoriaProfundidad.Simula}
  \leanok
  For a type $\sigma$, a map $\delta : \sigma \to \mathrm{Option}\ \sigma$, a state $s_0$ and a code
  $p$, the proposition $\mathrm{Simula}(\delta, s_0, p)$ says: for every $n$ such that the trace of
  $\delta$ from $s_0$ is \texttt{none} (halted) at step $n$, the \texttt{evaln} of $p$ on input $0$
  with budget $n$ halts. The trace is $n \mapsto (c \mapsto c.\mathrm{bind}\ \delta)^n(\mathrm{some}\
  s_0)$. The entry only says that the definition exists; no entry asserts it.
  \own{}
  \leanfile{Sagitta/Map/Informacion/MemoriaProfundidad.lean}
\end{definition}

\begin{theorem}[Memory bounds depth, given that the code is no slower than the system]
  \label{thm:information-memory-depth}
  \lean{Sagitta.Map.Informacion.MemoriaProfundidad.memoria_acota_profundidad}
  \leanok
  \uses{def:information-simula, def:information-prof}
  Let $\sigma$ be a finite type, $\delta : \sigma \to \mathrm{Option}\ \sigma$, $s_0 \in \sigma$ and
  $s \in \mathbb{N}$ with $|\sigma| \le 2^s$. Assume that the trace of $\delta$ from $s_0$ reaches
  \texttt{none} at some step. Let $p$ be a code and $x \in \mathbb{N}$ such that $p$ produces $x$,
  and assume $\mathrm{Simula}(\delta, s_0, p)$. Then $\mathrm{prof}(x) + 3 \le |p| + s$.
  \source{\cite{Bennett88, Antunes06}.}
  \fidelity{Not checked against the printed source.}
  \leanfile{Sagitta/Map/Informacion/MemoriaProfundidad.lean}
\end{theorem}

\begin{proof}
  \leanok
  A halting trace on at most $2^s$ states halts within $2^s$ steps, so $p$ runs within
  \texttt{evaln} budget $2^s$; then the bound $\mathrm{Kt}(x) \le |p| + \lfloor \log_2 2^s \rfloor$
  together with $\KC(x) \ge 3$ gives the claim.
\end{proof}

\section{The law of three memories, in a black-box model}
\label{sec:information-three-memories}

The model of this section has no machine. There is one hidden solution $y$ in a finite type $Y$ for
each instance. A verifier is a black box $O\,y : Q \to R$ that answers a query $q$ with $O\,y\,q$. A
solver is a query tree, and a family of query trees indexed by a program $p$ and a hint $c$ is a
solver with a description and a hint.

\begin{definition}[Query trees to a black-box verifier, and solving in $t$ queries]
  \label{def:information-query-tree}
  \lean{Sagitta.Map.Tiempo.LeyTresMemorias.ArbolV, Sagitta.Map.Tiempo.LeyTresMemorias.ArbolV.corre, Sagitta.Map.Tiempo.LeyTresMemorias.ArbolV.pasos, Sagitta.Map.Tiempo.LeyTresMemorias.igual, Sagitta.Map.Tiempo.LeyTresMemorias.Resuelve}
  \leanok
  A tree over types $Q, R, Y$ is either a leaf holding a candidate $y \in Y$, or a query $q \in Q$
  together with a subtree for each answer $r \in R$. Against a verifier $O : Q \to R$, the tree
  returns the candidate of the leaf that it reaches, and its time is the number of queries on that
  path. The yes-or-no verifier of the instance $y$ is $\mathrm{igual}(y) = (z \mapsto [z = y])$,
  for $Y$ with decidable equality. For $O : Y \to Q \to R$, trees $A : P \to C \to \text{tree}$ and
  $t \in \mathbb{N}$, the proposition $\mathrm{Resuelve}(O, A, t)$ says: for every $y \in Y$ there
  are $p \in P$ and $c \in C$ such that the tree $A\,p\,c$ against $O\,y$ returns $y$ in at most $t$
  queries.
  \own{}
  \leanfile{Sagitta/Map/LeyTresMemorias.lean}
\end{definition}

\begin{theorem}[The law of three memories, for a black-box verifier with answers of $b$ bits]
  \label{thm:information-three-memories}
  \lean{Sagitta.Map.Tiempo.LeyTresMemorias.ley, Sagitta.Map.Tiempo.LeyTresMemorias.ley_log}
  \leanok
  \uses{def:information-query-tree}
  Let $P, C, R, Y$ be finite types with decidable equality on $Y$, let $Q$ be a type, and let
  $L, s, b, t \in \mathbb{N}$ with $|P| \le 2^L$, $|C| \le 2^s$ and $|R| \le 2^b$. Let
  $O : Y \to Q \to R$ and $A : P \to C \to \text{tree}(Q, R, Y)$ with $\mathrm{Resuelve}(O, A, t)$.
  Then $|Y| \le 2^L \cdot 2^s \cdot 2^{b t}$, and also
  $\lceil \log_2 |Y| \rceil \le L + s + b t$.
  \own{}
  \leanfile{Sagitta/Map/LeyTresMemorias.lean}
\end{theorem}

\begin{proof}
  \leanok
  A tree that makes at most $t$ queries against answers from a set of size at most $2^b$ finds at
  most $(2^b)^t$ solutions, because a solution that it finds is determined by the sequence of
  answers it receives; the $|P||C|$ trees together cover $Y$. The second form takes logarithms.
\end{proof}

\begin{theorem}[Replacing the programs by one searching program multiplies the time by the number of programs, in a black-box model]
  \label{thm:information-levin-box}
  \lean{Sagitta.Map.Tiempo.LeyTresMemorias.levin_caja_ley}
  \leanok
  \uses{def:information-query-tree}
  Let $P, C, Y$ be finite types with decidable equality on $Y$, let $d \in Y$ and
  $L, s, t \in \mathbb{N}$ with $|P| \le 2^L$ and $|C| \le 2^s$. Let
  $A : P \to C \to \text{tree}(Y, \mathrm{Bool}, Y)$ with $\mathrm{Resuelve}(\mathrm{igual}, A, t)$.
  Then $|P|(t + 1) \le 2^L (t + 1)$ and
  $|Y| \le 2^0 \cdot 2^s \cdot (|P|(t + 1) + 1)$.
  \own{}
  \leanfile{Sagitta/Map/LeyTresMemorias.lean}
\end{theorem}

\begin{proof}
  \leanok
  The first inequality multiplies $|P| \le 2^L$ by $t + 1$. For the second, a universal tree that
  tries the programs one after another, each cut at $t$ queries, solves in $|P|(t+1)$ queries with a
  single program; the law for a yes-or-no verifier with one program applies to it.
\end{proof}

\section{The plane of costs}
\label{sec:information-cost-plane}

The first two statements are in the black-box model of the previous section, with the yes-or-no
verifier on $2^m$ candidates, $\mathrm{siNo}(m) : \mathrm{Fin}(2^m) \to \mathrm{Fin}(2^m) \to
\mathrm{Bool}$, which is $\mathrm{igual}$ on $\mathrm{Fin}(2^m)$. A point $(x, y) \in \mathbb{N}
\times \mathbb{N}$ is \emph{reached} when there are finite types $P$ and $C$, a family of trees
$A : P \to C \to \text{tree}$, and numbers $L, s, t$ with $|P| \le 2^L$, $|C| \le 2^s$,
$\mathrm{Resuelve}(\mathrm{siNo}(m), A, t)$, $L + s \le x$ and $\lceil \log_2 (t + 1) \rceil \le y$.
The \emph{region} is the set of reached points, and the \emph{frontier} is the set of points of the
region such that every point of the region that is at most it in both coordinates is equal to it.

\begin{theorem}[The region of the yes-or-no verifier is a half-plane, in a black-box model]
  \label{thm:information-half-plane}
  \lean{Sagitta.Map.Tiempo.PlanoCostos.region_es_semiplano}
  \leanok
  \uses{def:information-query-tree, thm:information-three-memories}
  For all $m, x, y \in \mathbb{N}$, the point $(x, y)$ is reached with the yes-or-no verifier on
  $2^m$ candidates if and only if $m \le x + y$.
  \own{}
  \leanfile{Sagitta/Map/PlanoCostos.lean}
\end{theorem}

\begin{proof}
  \leanok
  If the point is reached, the law of three memories in bits gives $m \le L + s + \lceil \log_2
  (t+1) \rceil \le x + y$. Conversely, an exact mixture of a hint of $m - k$ bits and a search of
  $2^k - 1$ queries reaches the point $(m - k, k)$ with $k = \min(y, m)$, and the region is closed
  upward.
\end{proof}

\begin{theorem}[The frontier of the yes-or-no verifier is the line $x + y = m$, in a black-box model]
  \label{thm:information-frontier-line}
  \lean{Sagitta.Map.Tiempo.PlanoCostos.frontera_es_recta}
  \leanok
  \uses{thm:information-half-plane}
  For every $m \in \mathbb{N}$ and every point $p = (p_1, p_2)$, $p$ is in the frontier of the
  region of the yes-or-no verifier on $2^m$ candidates if and only if $p_1 + p_2 = m$.
  \own{}
  \leanfile{Sagitta/Map/PlanoCostos.lean}
\end{theorem}

\begin{proof}
  \leanok
  By the half-plane description, a point with $p_1 + p_2 > m$ has a neighbour one step down in the
  region, and a point with $p_1 + p_2 = m$ is in the region with no point of the region below it
  other than itself.
\end{proof}

\begin{theorem}[A function of Blum speed-up, in the plane of costs]
  \label{thm:information-blum-plane}
  \lean{Sagitta.Map.Tiempo.PlanoCostos.blum_en_el_plano}
  \leanok
  % uses-lean: Sagitta.Map.Tiempo.o, Sagitta.Map.Tiempo.EsO
  \uses{thm:information-blum-sin-optimo, thm:information-blum-exists}
  There is $f : \mathbb{N} \to \mathbb{N}$ with $f(n) \le 1$ for all $n$, computable, such that
  \begin{enumerate}
    \item for every code $i$ that computes $f$ and every $M \in \mathbb{N}$ there is a code $j$
      that computes $f$ with $\mathrm{cost}(j, \cdot) \in o(\mathrm{cost}(i, \cdot))$ and
      $M < \mathrm{encode}(j)$;
    \item there is no code $i$ that computes $f$ and satisfies, for every code $j$ that computes
      $f$, $\mathrm{cost}(i, \cdot) = O(\mathrm{cost}(j, \cdot))$;
    \item there is no $M$ such that for every code $j$ that computes $f$ there is a code $i$ that
      computes $f$ with $\mathrm{encode}(i) \le M$ and $\mathrm{cost}(i, \cdot) =
      O(\mathrm{cost}(j, \cdot))$.
  \end{enumerate}
  Here $\mathrm{encode}$ is the natural number that Mathlib assigns to a code, and the cost is the
  cost of Blum.
  \own{}
  \leanfile{Sagitta/Map/PlanoCostos.lean}
\end{theorem}

\begin{proof}
  \leanok
  Take the function of Blum speed-up with the linear speed-up rate. From any program, a ladder of
  programs of $f$ descends, each with cost negligible against the previous one, and the codes along
  the ladder are pairwise different, hence not bounded; the second part is the absence of an
  optimal program, and the third repeats the argument for the programs with code at most $M$.
\end{proof}

\section{The closed universe}
\label{sec:information-closed-universe}

In the closed universe nothing enters for free. The machine is the evaluator of Mathlib codes with
a budget: $\mathrm{correr}(i, k, x)$ runs the code numbered $i$ on $x$ with \texttt{evaln} budget
$k$. The costs below are budgets of \texttt{evaln}, not machine steps, and the verifier is evaluated
at no charge in the first model and at a charge in the second one.

\begin{definition}[Closed universe, closed process, closed cost and $\mathrm{kt}$]
  \label{def:information-closed-universe}
  \lean{Sagitta.Map.Tiempo.UniversoCerrado.Universo, Sagitta.Map.Tiempo.UniversoCerrado.Proceso, Sagitta.Map.Tiempo.UniversoCerrado.Correcto, Sagitta.Map.Tiempo.UniversoCerrado.Resoluble, Sagitta.Map.Tiempo.UniversoCerrado.costoCerrado, Sagitta.Map.Tiempo.UniversoCerrado.kt}
  \leanok
  % uses-lean: Sagitta.Map.Levin.correr
  A universe $U$ is a verifier $V : \mathbb{N} \to \mathbb{N} \to \mathrm{Bool}$ and an input
  $x \in \mathbb{N}$. A process $p$ is a pair $(i, k)$ of a program number and a budget. The process
  is correct for $U$ when there is $y$ with $\mathrm{correr}(i, k, x) = \mathrm{some}\ y$ and
  $V(x, y) = \mathrm{true}$, and $U$ is solvable when some process is correct. The closed cost of
  $U$ is the infimum of $2^i k$ over the correct processes $(i, k)$ (it is $0$ when there is none),
  and $\mathrm{kt}(U)$ is the infimum of $i + \lceil \log_2 k \rceil$ over the correct processes.
  \own{}
  \leanfile{Sagitta/Map/UniversoCerrado.lean}
\end{definition}

\begin{theorem}[Levin's search costs less than eight times the closed cost, with a free verifier]
  \label{thm:information-closed-law}
  \lean{Sagitta.Map.Tiempo.UniversoCerrado.ley}
  \leanok
  % uses-lean: Sagitta.Map.Levin.buscar, Sagitta.Map.Tiempo.CorteRacional.costoLevin
  \uses{def:information-closed-universe}
  Let $U$ be a solvable closed universe. Then there are $T$ and $y$ such that the universal search
  of Levin on $(U.V, U.x)$ up to round $T$ returns $y$ and $U.V(U.x, y) = \mathrm{true}$, and the
  cost of the search on $(U.V, U.x)$ is less than $8$ times the closed cost of $U$. In the search,
  round $t$ runs the programs $0, \dots, t$ with budget $2^{t - i}$ for program $i$ and keeps the
  first output that the verifier accepts; the cost is the total budget handed out up to the first
  round that returns an answer.
  \own{}
  \leanfile{Sagitta/Map/UniversoCerrado.lean}
\end{theorem}

\begin{proof}
  \leanok
  Take a correct process that attains the closed cost. Its budget exceeds the input, so it is at
  least $1$, and the bound of Levin for one program, namely that the search costs less than
  $8 \cdot 2^i k$, applies to it.
\end{proof}

\begin{theorem}[The closed cost is attained, is below that of every correct process, and exceeds the input]
  \label{thm:information-closed-lower}
  \lean{Sagitta.Map.Tiempo.UniversoCerrado.cota_inferior}
  \leanok
  \uses{def:information-closed-universe}
  Let $U$ be a solvable closed universe. Then there is a correct process $p = (i, k)$ with closed
  cost of $U$ equal to $2^i k$; for every correct process $(i, k)$ the closed cost of $U$ is at
  most $2^i k$; and $U.x$ is less than the closed cost of $U$.
  \own{}
  \leanfile{Sagitta/Map/UniversoCerrado.lean}
\end{theorem}

\begin{proof}
  \leanok
  The set of costs of correct processes is a nonempty set of natural numbers, so its infimum is
  attained. A process that returns a value has a budget larger than its input, and the factor
  $2^i$ is at least $1$.
\end{proof}

\begin{theorem}[$\mathrm{kt}$ of a closed universe is attained]
  \label{thm:information-kt-attained}
  \lean{Sagitta.Map.Tiempo.UniversoCerrado.kt_alcanzado}
  \leanok
  \uses{def:information-closed-universe}
  Let $U$ be a solvable closed universe. Then there is a correct process $(i, k)$ with
  $\mathrm{kt}(U) = i + \lceil \log_2 k \rceil$.
  \own{}
  \leanfile{Sagitta/Map/UniversoCerrado.lean}
\end{theorem}

\begin{proof}
  \leanok
  The infimum of a nonempty set of natural numbers is a member of the set.
\end{proof}

\begin{definition}[Closed universe with a paid verifier]
  \label{def:information-paid-universe}
  \lean{Sagitta.Map.Tiempo.VerificadorPagado.UniversoPagado, Sagitta.Map.Tiempo.VerificadorPagado.ProcesoPagado, Sagitta.Map.Tiempo.VerificadorPagado.CorrectoPagado, Sagitta.Map.Tiempo.VerificadorPagado.ResolublePagado, Sagitta.Map.Tiempo.VerificadorPagado.costoCerradoPagado, Sagitta.Map.Tiempo.VerificadorPagado.buscarPagado, Sagitta.Map.Tiempo.VerificadorPagado.costoLevinPagado}
  \leanok
  A paid universe $U$ is a code $v$ and an input $x$; a paid process is a triple $(i, k, K)$. It
  is correct when there is $y$ with $\mathrm{correr}(i, k, x) = \mathrm{some}\ y$ and the
  \texttt{evaln} of $v$ on $\mathrm{pair}(x, y)$ with budget $K$ is $\mathrm{some}\ 0$; the universe
  is solvable when some paid process is correct. The paid cost of a process is $2^i (k + K)$, and
  the paid closed cost of $U$ is the infimum of the paid costs of the correct processes. The paid
  search in round $t$ runs program $i$ with budget $2^{t-i}$ and keeps its output $y$ only when the
  \texttt{evaln} of $v$ on $\mathrm{pair}(x, y)$ with the same budget gives $\mathrm{some}\ 0$;
  $\mathrm{buscarPagado}(v, x, T)$ is the first answer kept in rounds $0, \dots, T$, and
  $\mathrm{costoLevinPagado}(v, x)$ is twice the budget handed out up to the first round that
  returns an answer.
  \own{}
  \leanfile{Sagitta/Map/VerificadorPagado.lean}
\end{definition}

\begin{theorem}[Levin's search that pays for verification costs less than sixteen times the paid closed cost]
  \label{thm:information-paid-law}
  \lean{Sagitta.Map.Tiempo.VerificadorPagado.ley_pagada}
  \leanok
  \uses{def:information-paid-universe}
  Let $U = (v, x)$ be a solvable paid universe. Then there are $T$ and $y$ with
  $\mathrm{buscarPagado}(v, x, T) = \mathrm{some}\ y$ and $0$ is among the values of $v$ on
  $\mathrm{pair}(x, y)$, and $\mathrm{costoLevinPagado}(v, x) < 16 \cdot$ the paid closed cost of $U$.
  \own{}
  \leanfile{Sagitta/Map/VerificadorPagado.lean}
\end{theorem}

\begin{proof}
  \leanok
  As for the free verifier: take a correct paid process that attains the paid closed cost and apply
  the bound for one process, which now has the factor $16$ because each program receives a budget to
  run and the same budget to verify.
\end{proof}

\begin{theorem}[The minimum survives exactly when verifying is cheaper than every solver, with a paid verifier]
  \label{thm:information-paid-survives}
  \lean{Sagitta.Map.Tiempo.VerificadorPagado.sobrevive_iff}
  \leanok
  \uses{def:information-paid-universe}
  Let $v$ be a code and $f : \mathbb{N} \to \mathbb{N}$ such that $v$ verifies $f$, that is, $v$ is
  total and $0$ is a value of $v$ on $\mathrm{pair}(x, y)$ if and only if $y = f(x)$. Let $g :
  \mathbb{N} \to \mathbb{N}$, and let $j_0$ be a code that computes $f$. Write $t_V(x) =
  \mathrm{cost}(v, \mathrm{pair}(x, f(x)))$. Then the following are equivalent: for every code $j$
  that computes $f$, $x \mapsto$ (paid closed cost of $(v, x)$) is $O(x \mapsto \mathrm{cost}(j, x)
  + g(x))$; and for every code $j$ that computes $f$, $t_V$ is $O(x \mapsto \mathrm{cost}(j, x) +
  g(x))$.
  \own{}
  \leanfile{Sagitta/Map/VerificadorPagado.lean}
\end{theorem}

\begin{proof}
  \leanok
  The paid closed cost is at least $t_V$, and for a code $j$ that computes $f$ it is at most
  $2^j (\mathrm{cost}(j, x) + t_V(x))$; each direction composes one of these bounds with the
  hypothesis.
\end{proof}

\section{Paying once or at each instance}
\label{sec:information-two-accounts}

The first two statements are in the black-box model of Section~\ref{sec:information-three-memories}
with the yes-or-no verifier, for a family of sizes $m \in \mathbb{N}$: for each $m$ a finite type
$Y_m$ of hidden solutions with $|Y_m| = 2^m$, finite types $P_m$ and $C_m$ of programs and hints, and
a family of trees $A_m : P_m \to C_m \to \text{tree}(Y_m, \mathrm{Bool}, Y_m)$, with bounds
$|P_m| \le 2^{L_m}$ and $|C_m| \le 2^{s_m}$ and
$\mathrm{Resuelve}(\mathrm{igual}, A_m, t_m)$ for every $m$.

\begin{theorem}[The price is paid: description, hint and time add up to the size, in a black-box model]
  \label{thm:information-price-paid}
  \lean{Sagitta.Map.Tiempo.DosCuentas.precio_se_paga}
  \leanok
  \uses{def:information-query-tree, thm:information-three-memories}
  Under the hypotheses above, with $L, s, t : \mathbb{N} \to \mathbb{N}$, for every $m \in
  \mathbb{N}$, $m \le L_m + (s_m + \lceil \log_2 (t_m + 1) \rceil)$.
  \own{}
  \leanfile{Sagitta/Map/DosCuentas.lean}
\end{theorem}

\begin{proof}
  \leanok
  The law of three memories for the yes-or-no verifier in bits, at size $m$, with
  $\lceil \log_2 2^m \rceil = m$.
\end{proof}

\begin{theorem}[Paying once does not suffice, in a black-box model]
  \label{thm:information-once-not-enough}
  \lean{Sagitta.Map.Tiempo.DosCuentas.una_vez_no_alcanza}
  \leanok
  \uses{thm:information-price-paid}
  Under the same hypotheses, let $L_0 \in \mathbb{N}$ with $L_m \le L_0$ for every $m$. Then for
  every $m$, $m - L_0 \le s_m + \lceil \log_2 (t_m + 1) \rceil$ (subtraction of natural numbers),
  and if $L_0 < m$ then $1 \le s_m + \lceil \log_2 (t_m + 1) \rceil$.
  \own{}
  \leanfile{Sagitta/Map/DosCuentas.lean}
\end{theorem}

\begin{proof}
  \leanok
  Theorem~\ref{thm:information-price-paid} together with $L_m \le L_0$, by linear arithmetic.
\end{proof}

\begin{definition}[Machine $\mathrm{kt}$ is logarithmic on a set, on cslib machines]
  \label{def:information-ktmt-log}
  \lean{Sagitta.Map.Tiempo.LevinMaquinas.ResuelveMT, Sagitta.Map.Tiempo.LevinMaquinas.ResolubleMT, Sagitta.Map.Tiempo.LevinMaquinas.ktMT, Sagitta.Map.Tiempo.LevinMaquinas.KtMTLog}
  \leanok
  % uses-lean: Cslib.Complexity.SAT
  \uses{def:information-machine-k}
  Let $\alpha$ be a type with an embedding $\mathrm{enc} : \alpha \hookrightarrow \{0,1\}^*$ and
  $V : \alpha \times \{0,1\}^* \to \mathrm{Bool}$. The machine numbered $i$ \emph{solves} $a$ with
  clock $t$ when $\texttt{correrMT}(i, \mathrm{enc}(a), t) = \mathrm{some}\ y$ and
  $V(a, y) = \mathrm{true}$ for some $y$; $a$ is \emph{solvable} when some $i$ and $t$ do so. The
  number $\mathrm{ktMT}(a)$ is the infimum of $i + \lceil \log_2 t \rceil$ over the pairs $(i, t)$
  such that machine $i$ solves $a$ with clock $t$ ($0$ if there is none). For $A \subseteq \alpha$,
  $\mathrm{KtMTLog}(A)$ says: there is $c$ such that every $a \in A$ is solvable and
  $\mathrm{ktMT}(a) \le c \cdot \lfloor \log_2 (|\mathrm{enc}(a)| + 2) \rfloor$.
  \own{}
  \leanfile{Sagitta/Map/LevinMaquinas.lean}
\end{definition}

\begin{definition}[A family of satisfiable formulas of the second case]
  \label{def:information-case-two}
  \lean{Sagitta.Map.Tiempo.DosCasos.CasoDos}
  \leanok
  \uses{def:information-ktmt-log}
  For a set $A$ of CNF formulas over the natural numbers, $\mathrm{CasoDos}(A)$ says: $A$ is a
  subset of $\SAT$, the set of satisfiable formulas, and it is not the case that
  $\mathrm{KtMTLog}(A)$ holds for the encoding of formulas and the verifier of $\SAT$ of cslib.
  \own{}
  \leanfile{Sagitta/Map/DosCasos.lean}
\end{definition}

\begin{theorem}[A family of the second case exists if and only if machine $\mathrm{kt}$ of $\SAT$ is not logarithmic, on cslib machines]
  \label{thm:information-case-two-iff}
  \lean{Sagitta.Map.Tiempo.DosCasos.caso_dos_iff_no_kt_log}
  \leanok
  \uses{def:information-case-two}
  There is a set $A$ with $\mathrm{CasoDos}(A)$ if and only if $\mathrm{KtMTLog}$ does not hold for
  $\SAT$, with the encoding of formulas and the verifier of $\SAT$ of cslib.
  \own{}
  \leanfile{Sagitta/Map/DosCasos.lean}
\end{theorem}

\begin{proof}
  \leanok
  If $A \subseteq \SAT$ and $\mathrm{KtMTLog}$ held for $\SAT$, it would hold for the subset $A$.
  Conversely take $A = \SAT$.
\end{proof}

\section{The creator and the solver}
\label{sec:information-creator}

A creator draws $r$ random bits, computes a hidden solution from them, and declares a number for
the work of the computation. The number is declared; no theorem measures it against a model of
computation. The solver sees the instance only through the yes-or-no verifier.

\begin{definition}[Creator, and solving the family of a creator, in a black-box model]
  \label{def:information-creator}
  \lean{Sagitta.Map.Tiempo.HuevoGallina.Creador, Sagitta.Map.Tiempo.HuevoGallina.Creador.imagen, Sagitta.Map.Tiempo.HuevoGallina.ResuelveFamilia, Sagitta.Map.Tiempo.HuevoGallina.constante}
  \leanok
  \uses{def:information-query-tree}
  A creator $K$ over $Y$ is a number $r$ of random bits, a map $g : \mathrm{Fin}(2^r) \to Y$ from the
  seed to the hidden solution, and a number $\mathrm{trabajo}$. Its image is the set of values of
  $g$. For $g : \mathrm{Fin}(2^r) \to Y$, a family $A : C \to \text{tree}(Y, \mathrm{Bool}, Y)$ and
  $t \in \mathbb{N}$, $\mathrm{ResuelveFamilia}(g, A, t)$ says: for every seed $\sigma$ there is a
  hint $c$ such that $A\,c$ against $\mathrm{igual}(g(\sigma))$ returns $g(\sigma)$ in at most $t$
  queries. The constant creator $\mathrm{constante}(r, W)$ is the creator over the one-point type
  with $r$ random bits, the constant map, and $\mathrm{trabajo} = W$.
  \own{}
  \leanfile{Sagitta/Map/HuevoGallina.lean}
\end{definition}

\begin{theorem}[Information, not work, is what the solver pays: a constant creator, in a black-box model]
  \label{thm:information-not-work}
  \lean{Sagitta.Map.Tiempo.HuevoGallina.informacion_no_trabajo}
  \leanok
  \uses{def:information-creator}
  For all $r, W \in \mathbb{N}$: the creator $\mathrm{constante}(r, W)$ has $r$ random bits and
  work $W$, its image has exactly one element, and $\mathrm{ResuelveFamilia}$ holds for its map, the
  family that is the leaf of the unique point at every hint, and $t = 0$.
  \own{}
  \leanfile{Sagitta/Map/HuevoGallina.lean}
\end{theorem}

\begin{proof}
  \leanok
  Unfold the definitions: the image of a constant map from a nonempty type is a singleton, and a
  leaf makes no query.
\end{proof}

\begin{theorem}[Asymmetry between creating and solving, in a black-box model]
  \label{thm:information-asymmetry}
  \lean{Sagitta.Map.Tiempo.HuevoGallina.asimetria}
  \leanok
  \uses{def:information-creator}
  For every $r \in \mathbb{N}$ there is a creator $K$ over $\mathrm{Fin}(2^r)$ with $K.r = r$ and
  $K.\mathrm{trabajo} = 1$ such that
  \begin{enumerate}
    \item for every family $A : \mathrm{Unit} \to \text{tree}$ and every $t$, if
      $\mathrm{ResuelveFamilia}(K.g, A, t)$ then $2^r \le t + 1$;
    \item for all $s, t \in \mathbb{N}$, every finite type $C$ with $|C| \le 2^s$ and every
      $A : C \to \text{tree}$, if $\mathrm{ResuelveFamilia}(K.g, A, t)$ then $2^r \le 2^s (t+1)$;
    \item $\mathrm{ResuelveFamilia}(K.g, \sigma \mapsto \mathrm{leaf}(K.g(\sigma)), 0)$ holds.
  \end{enumerate}
  Here the trees are over $\mathrm{Fin}(2^r)$ with answers in $\mathrm{Bool}$ and candidates in
  $\mathrm{Fin}(2^r)$.
  \own{}
  \leanfile{Sagitta/Map/HuevoGallina.lean}
\end{theorem}

\begin{proof}
  \leanok
  Take the identity creator on $\mathrm{Fin}(2^r)$. Parts~1 and~2 are the bound for the image of an
  injective map, since a hint reaches at most $t + 1$ solutions; part~3 is the seed as the hint.
\end{proof}

\section{Information is paid once}
\label{sec:information-amortization}

\begin{definition}[Total cost of a table shared by $M$ queries]
  \label{def:information-total}
  \lean{Sagitta.Map.Tiempo.LeyInformacion.total}
  \leanok
  For $T, M, r \in \mathbb{N}$, $\mathrm{total}(T, M, r) = T + M r$: a table of declared cost $T$,
  paid once, and $M$ queries of declared cost $r$ each. The costs are declared numbers, not
  measured in a model of computation.
  \own{}
  \leanfile{Sagitta/Map/LeyInformacion.lean}
\end{definition}

\begin{theorem}[Someone paid for the table]
  \label{thm:information-someone-paid}
  \lean{Sagitta.Map.Tiempo.LeyInformacion.alguien_pago}
  \leanok
  \uses{def:information-total}
  For all $T, M, r \in \mathbb{N}$, $T \le \mathrm{total}(T, M, r)$.
  \own{}
  \leanfile{Sagitta/Map/LeyInformacion.lean}
\end{theorem}

\begin{proof}
  \leanok
  The total is $T$ plus a natural number.
\end{proof}

\begin{theorem}[Sharing the table is cheaper than paying for it at each query]
  \label{thm:information-amortized}
  \lean{Sagitta.Map.Tiempo.LeyInformacion.total_amortizado}
  \leanok
  \uses{def:information-total}
  For all $T, M, r \in \mathbb{N}$ with $1 \le M$, $\mathrm{total}(T, M, r) \le M (T + r)$.
  \own{}
  \leanfile{Sagitta/Map/LeyInformacion.lean}
\end{theorem}

\begin{proof}
  \leanok
  Expand $M (T + r) = M T + M r$ and use $T \le M T$.
\end{proof}

\section{$\SAT$ in the closed universe}
\label{sec:information-sat-closed}

Formulas are lists of clauses, which is the type of CNF formulas over the natural numbers of
cslib-frontier. The verifier $\mathrm{vSAT}(x, y)$ decodes $x$ as a formula $F$ and $y$ as a list of
Booleans, read as an assignment that is false outside the list, and returns the value of $F$ under
it; it returns false if either does not decode. The closed universe of a formula $F$ is
$\mathrm{uSAT}(F) = (\mathrm{vSAT}, \mathrm{encode}(F))$ in the sense of
Section~\ref{sec:information-closed-universe}, and $\mathrm{ktSAT}(F) = \mathrm{kt}(\mathrm{uSAT}(F))$.

\begin{definition}[$\SAT$ is easy in the closed universe, and $\mathrm{kt}$ is logarithmic]
  \label{def:information-sat-easy}
  \lean{Sagitta.Map.Tiempo.KtSAT.vSAT, Sagitta.Map.Tiempo.KtSAT.Satisfacible, Sagitta.Map.Tiempo.KtSAT.uSAT, Sagitta.Map.Tiempo.KtSAT.ktSAT, Sagitta.Map.Tiempo.KtSAT.FacilCerrado, Sagitta.Map.Tiempo.KtSAT.KtLogaritmico}
  \leanok
  \uses{def:information-closed-universe}
  $\mathrm{FacilCerrado}$ is the proposition: there are $C, d \in \mathbb{N}$ such that for every
  satisfiable formula $F$, the closed cost of $\mathrm{uSAT}(F)$ is at most
  $C (\mathrm{encode}(F) + 2)^d$. $\mathrm{KtLogaritmico}$ is the proposition: there is
  $c \in \mathbb{N}$ such that for every satisfiable formula $F$,
  $\mathrm{ktSAT}(F) \le c \lfloor \log_2 (\mathrm{encode}(F) + 2) \rfloor$. The entry only says that
  the definitions exist; no entry asserts either proposition.
  \own{}
  \leanfile{Sagitta/Map/KtSAT.lean}
\end{definition}

\begin{theorem}[$\SAT$ is easy in the closed universe if and only if $\mathrm{kt}$ is logarithmic]
  \label{thm:information-sat-easy-iff-kt}
  \lean{Sagitta.Map.Tiempo.KtSAT.facil_iff_kt}
  \leanok
  \uses{def:information-sat-easy}
  $\mathrm{FacilCerrado}$ holds if and only if $\mathrm{KtLogaritmico}$ holds.
  \own{}
  \leanfile{Sagitta/Map/KtSAT.lean}
\end{theorem}

\begin{proof}
  \leanok
  From a polynomial bound $C n^d$ on the closed cost, the process that attains $\mathrm{kt}$ is
  below it in bits, with a constant $c$ made from $C$ and $d$; from $c$, the closed cost is at most
  $2^{\mathrm{kt}}$, a polynomial of degree $c$.
\end{proof}

\begin{theorem}[$\mathrm{kt}$ of a satisfiable formula is the logarithm of the closed cost, up to a constant]
  \label{thm:information-sat-law}
  \lean{Sagitta.Map.Tiempo.KtSAT.ley_sat}
  \leanok
  \uses{def:information-sat-easy, thm:information-closed-law}
  Let $F$ be a satisfiable formula. Write $c_1 = \lfloor \log_2 (\text{closed cost of }
  \mathrm{uSAT}(F)) \rfloor$ and $c_2 = \lfloor \log_2 (\mathrm{costoLevin}(\mathrm{vSAT},
  \mathrm{encode}(F))) \rfloor$. Then $c_1 \le \mathrm{ktSAT}(F)$, $\mathrm{ktSAT}(F) \le c_1 + 1$,
  $\mathrm{ktSAT}(F) \le c_2$ and $c_2 < \mathrm{ktSAT}(F) + 3$.
  \own{}
  \leanfile{Sagitta/Map/KtSAT.lean}
\end{theorem}

\begin{proof}
  \leanok
  A satisfiable formula has a solvable closed universe. The four inequalities are the four
  conversions to bits that hold in every solvable closed universe, between $\mathrm{kt}$, the closed
  cost and the cost of Levin's search.
\end{proof}

\begin{definition}[A model of a proof system]
  \label{def:information-proof-system}
  \lean{Sagitta.Map.Tiempo.KtPruebaFija.SistemaPruebas, Sagitta.Map.Tiempo.KtPruebaFija.vPrueba, Sagitta.Map.Tiempo.KtPruebaFija.uPrueba, Sagitta.Map.Tiempo.KtPruebaFija.indiceFijo}
  \leanok
  \uses{def:information-closed-universe}
  A proof system $S$ is a function $\mathrm{verifica} : \mathbb{N} \to \mathbb{N} \to \mathrm{Bool}$
  and a function $\mathrm{largo} : \mathbb{N} \to \mathbb{N}$, with the reading that
  $\mathrm{verifica}(T, \pi) = \mathrm{true}$ means that $\pi$ proves $T$. The verifier
  $\mathrm{vPrueba}(S)(x, \pi)$ is true when $\mathrm{verifica}(T', \pi)$ is true and
  $\mathrm{largo}(\pi) \le n'$, where $(T', n') = \mathrm{unpair}(x)$. The universe
  $\mathrm{uPrueba}(S, T, n)$ is $(\mathrm{vPrueba}(S), \mathrm{pair}(T, n))$. For $y_0 \in
  \mathbb{N}$, $\mathrm{indiceFijo}(y_0)$ is the number of the code of Mathlib that, on every input,
  returns $y_0$ by its tree.
  \own{}
  \leanfile{Sagitta/Map/KtPruebaFija.lean}
\end{definition}

\begin{theorem}[A fixed theorem: its proof is paid once, in the closed universe]
  \label{thm:information-fixed-theorem}
  \lean{Sagitta.Map.Tiempo.KtPruebaFija.teorema_fijo}
  \leanok
  \uses{def:information-proof-system, thm:information-kt-attained}
  Let $S$ be a proof system and $T, \pi \in \mathbb{N}$ with $S.\mathrm{verifica}(T, \pi) =
  \mathrm{true}$, and let $n \in \mathbb{N}$ with $S.\mathrm{largo}(\pi) \le n$. Then the universe
  $\mathrm{uPrueba}(S, T, n)$ is solvable, $\lfloor \log_2 \mathrm{pair}(T, n) \rfloor \le
  \mathrm{kt}(\mathrm{uPrueba}(S, T, n))$, and $\mathrm{kt}(\mathrm{uPrueba}(S, T, n)) \le
  \mathrm{indiceFijo}(\pi) + \lceil \log_2 (\mathrm{pair}(T, n) + 1) \rceil$ and
  $\mathrm{kt}(\mathrm{uPrueba}(S, T, n)) \le \mathrm{indiceFijo}(\pi) + 2 \lceil \log_2 (\max(T, n) + 1) \rceil$.
  \own{}
  \leanfile{Sagitta/Map/KtPruebaFija.lean}
\end{theorem}

\begin{proof}
  \leanok
  The process made of the fixed code that writes $\pi$ and the budget $\mathrm{pair}(T, n) + 1$ is
  correct; this gives solvability and the upper bounds, and the lower bound is that a process that
  returns a value reads its input.
\end{proof}

\section{The third account: justifying a ``no''}
\label{sec:information-third-account}

In the closed universe of Section~\ref{sec:information-closed-universe} the verifier checks every
``yes'' with its certificate. A ``no'' has to be justified by a refutation that the universe
checks. Here the refutations are those of resolution in bounded space. A refutation in space $s$
is a list of moves, each of which loads a clause of the formula, adds the resolvent of two clauses
that are in memory, or deletes a clause, such that memory never holds more than $s$ clauses and the
run ends with the empty clause in memory. Its length counts loads and resolutions. The number
$\mathrm{largoMinimo}(F, s) \in \mathbb{N} \cup \{\infty\}$ is the least length of a refutation of
$F$ in space $s$.

\begin{definition}[Verdicts, their price, and the third account]
  \label{def:information-third-account}
  \lean{Sagitta.Map.Tiempo.TerceraCuenta.Veredicto, Sagitta.Map.Tiempo.TerceraCuenta.aceptaJ, Sagitta.Map.Tiempo.TerceraCuenta.precioJ, Sagitta.Map.Tiempo.TerceraCuenta.terceraCuenta, Sagitta.Map.Tiempo.CajaGris.largoMinimo}
  \leanok
  A verdict on a formula $F$ is either $\mathrm{si}(\sigma)$, with a list $\sigma$ of Booleans read
  as an assignment that is false outside the list, or $\mathrm{no}(s, ms)$, with a space $s$ and a
  list $ms$ of moves. The universe accepts $\mathrm{si}(\sigma)$ when $F$ evaluates to true under
  $\sigma$, and $\mathrm{no}(s, ms)$ when $ms$ is a legal run in space $s$ that ends with the empty
  clause. The price of $\mathrm{si}(\sigma)$ is the number of literals read by the evaluation of $F$
  under $\sigma$, clause by clause until the first false one, and the price of $\mathrm{no}(s, ms)$
  is the length of $ms$. The third account $\mathrm{terceraCuenta}(F) \in \mathbb{N} \cup \{\infty\}$
  is the infimum of the prices of the accepted verdicts.
  \own{}
  \leanfile{Sagitta/Map/TerceraCuenta.lean}
\end{definition}

\begin{theorem}[The third account of an unsatisfiable formula is its shortest refutation in the best space]
  \label{thm:information-third-account-length}
  \lean{Sagitta.Map.Tiempo.TerceraCuenta.terceraCuenta_eq_largo}
  \leanok
  \uses{def:information-third-account}
  Let $F$ be a list of clauses that is not satisfiable. Then
  $\mathrm{terceraCuenta}(F) = \inf_s \mathrm{largoMinimo}(F, s)$.
  \own{}
  \leanfile{Sagitta/Map/TerceraCuenta.lean}
\end{theorem}

\begin{proof}
  \leanok
  A refutation in space $s$ is an accepted verdict $\mathrm{no}(s, ms)$ of the same price, so the
  left side is at most the right side; an accepted verdict $\mathrm{si}(\sigma)$ is impossible for
  an unsatisfiable formula, and an accepted $\mathrm{no}(s, ms)$ is a refutation in space $s$.
\end{proof}

\begin{theorem}[The third account of the pigeonhole formula is at least $2^{\lfloor n/24 \rfloor}$, in resolution]
  \label{thm:information-third-account-pigeonhole}
  \lean{Sagitta.Map.Tiempo.TerceraCuenta.terceraCuenta_palomar}
  \leanok
  \uses{def:information-third-account}
  Let $\mathrm{formulaPal}(n)$ be the pigeonhole formula with $n + 1$ pigeons and $n$ holes, with
  its variables renamed by a fixed order. For every $n \in \mathbb{N}$,
  $2^{\lfloor n/24 \rfloor} \le \mathrm{terceraCuenta}(\mathrm{formulaPal}(n))$, as elements of
  $\mathbb{N} \cup \{\infty\}$.
  \own{}
  \leanfile{Sagitta/Map/TerceraCuenta.lean}
\end{theorem}

\begin{proof}
  \leanok
  Every accepted verdict on the pigeonhole formula is a $\mathrm{no}$, since the formula has no
  satisfying assignment, and every accepted refutation, in any space, has length at least
  $2^{\lfloor n/24 \rfloor}$ by the lower bound for resolution proved in cslib-frontier.
\end{proof}

\section{Conservation of work, in a black-box model}
\label{sec:information-work}

In this section the object is a finite type $\alpha$ of options and an oracle $f : \alpha \to
\mathrm{Bool}$ that is seen only by asking.

\begin{definition}[Query algorithms, fillings and their chain, in a black-box model]
  \label{def:information-filling}
  \lean{Sagitta.Map.Tiempo.Pregunta, Sagitta.Map.Tiempo.Pregunta.corre, Sagitta.Map.Tiempo.Pregunta.vistos, Sagitta.Map.Tiempo.Pregunta.costo, Sagitta.Map.Tiempo.DecideExiste, Sagitta.Map.Tiempo.ConservacionTrabajo.Llenado, Sagitta.Map.Tiempo.ConservacionTrabajo.Llenado.tablaDe, Sagitta.Map.Tiempo.ConservacionTrabajo.Llenado.costoLl, Sagitta.Map.Tiempo.ConservacionTrabajo.encadenar}
  \leanok
  A query algorithm over $\alpha$ is either an answer in $\mathrm{Bool}$, or a query of an option
  $i \in \alpha$ together with a continuation for each of the two answers. Its cost on $f$ is the
  number of queries on the path that $f$ selects. It decides existence when, for every $f$, it
  answers $\mathrm{true}$ if and only if $f(i) = \mathrm{true}$ for some $i$. A filling over
  $\alpha$ with tables in a type $\tau$ is a tree of the same shape whose leaves hold a table
  $T \in \tau$ instead of an answer; $\mathrm{tablaDe}(L, f)$ is the table at the leaf that $f$
  selects and $\mathrm{costoLl}(L, f)$ is the number of queries on that path. Given $Q : \tau \to$
  query algorithm, $\mathrm{encadenar}(L, Q)$ is the query algorithm that runs the filling $L$ and
  then $Q(T)$ with the table $T$ found.
  \own{}
  \leanfile{Sagitta/Map/ConservacionTrabajo.lean}
\end{definition}

\begin{theorem}[Work is conserved between filling a table and querying it, in a black-box model]
  \label{thm:information-work-conserved}
  \lean{Sagitta.Map.Tiempo.ConservacionTrabajo.trabajo_se_conserva}
  \leanok
  \uses{def:information-filling}
  Let $\alpha$ and $\tau$ be types with $\alpha$ finite, $L$ a filling over $\alpha$ with tables in
  $\tau$ and $Q : \tau \to$ query algorithm. Assume that $\mathrm{encadenar}(L, Q)$ decides
  existence. Then $|\alpha| \le \mathrm{costoLl}(L, \mathrm{nula}) + \mathrm{costo}(Q(T_0),
  \mathrm{nula})$, where $\mathrm{nula}$ is the oracle that is false on every option and $T_0 =
  \mathrm{tablaDe}(L, \mathrm{nula})$.
  \own{}
  \leanfile{Sagitta/Map/ConservacionTrabajo.lean}
\end{theorem}

\begin{proof}
  \leanok
  The cost of the chain is the sum of the costs of the two phases. An algorithm that decides
  existence must look at every option when the oracle is $\mathrm{nula}$, because an option left
  unseen could be changed to $\mathrm{true}$ without the algorithm noticing.
\end{proof}

\begin{definition}[A table as a bounded-memory detector in backtracking]
  \label{def:information-table-detector}
  \lean{Sagitta.Map.Tiempo.ConservacionTrabajo.tablaLlena, Sagitta.Map.Tiempo.ConservacionTrabajo.memoriaTabla, Sagitta.Map.Tiempo.Vivo, Sagitta.Map.Tiempo.SanoN, Sagitta.Map.Tiempo.DemoraN, Sagitta.Map.Tiempo.buscar, Sagitta.Map.Tiempo.nodos}
  \leanok
  Let $\mathrm{sol} : \mathrm{List}(\mathrm{Bool}) \to \mathrm{Bool}$ be a predicate on lists of
  bits. A prefix $p$ is alive with $k$ bits to fix when some list $q$ of length $k$ makes
  $\mathrm{sol}(p \mathbin{+\!\!+} q)$ true. A detector $\det : \mathrm{List}(\mathrm{Bool}) \to
  \mathrm{Bool}$ is sound with horizon $N$ when for every prefix $p$ with $|p| \le N$, $\det(p) =
  \mathrm{true}$ implies that $p$ is not alive with $N - |p|$ bits to fix. It has delay $d$ with
  horizon $N$ when for every prefix $p$ with $|p| + d \le N$ that is not alive with $N - |p|$ bits
  to fix, $\det(p \mathbin{+\!\!+} q) = \mathrm{true}$ for every $q$ of length $d$. Backtracking
  with $\det$ at a prefix $p$ with $k$ bits to fix returns at once if $\det(p)$ holds, and
  otherwise tries the extension by $\mathrm{false}$ and, if that fails, by $\mathrm{true}$;
  $\mathrm{nodos}$ is the number of nodes it visits. For $N \in \mathbb{N}$,
  $\mathrm{tablaLlena}(\mathrm{sol}, N)$ is the table that says, for each prefix $p$, that $p$ is
  not alive with $N - |p|$ bits to fix, and $\mathrm{memoriaTabla}(T)$ is the detector whose state
  is the prefix read and whose alarm is the table $T$.
  \own{}
  \leanfile{Sagitta/Map/ConservacionTrabajo.lean}
\end{definition}

\begin{theorem}[A full table makes backtracking linear, with no delay]
  \label{thm:information-perfect-table}
  \lean{Sagitta.Map.Tiempo.ConservacionTrabajo.tabla_perfecta}
  \leanok
  \uses{def:information-table-detector}
  For every $\mathrm{sol} : \mathrm{List}(\mathrm{Bool}) \to \mathrm{Bool}$ and $N \in \mathbb{N}$,
  the detector $\det = \mathrm{memoriaTabla}(\mathrm{tablaLlena}(\mathrm{sol}, N))$ is sound with
  horizon $N$, has delay $0$ with horizon $N$, and backtracking with $\det$ from the empty prefix
  with $N$ bits to fix visits at most $2N + 1$ nodes.
  \own{}
  \leanfile{Sagitta/Map/ConservacionTrabajo.lean}
\end{theorem}

\begin{proof}
  \leanok
  On prefixes of length at most $N$ the detector answers exactly ``not alive'', which gives
  soundness and delay $0$; with a perfect detector the search visits one path, plus one failed
  child per level.
\end{proof}

\begin{theorem}[Work is conserved against randomness, in a black-box model]
  \label{thm:information-work-random}
  \lean{Sagitta.Map.Tiempo.ConservacionDual.trabajo_se_conserva_azar}
  \leanok
  \uses{def:information-filling, thm:information-work-conserved}
  Let $\alpha$ be a finite type with decidable equality and at least one element, let $\iota$ and
  $\tau$ be types, $S$ a nonempty finite set of seeds in $\iota$, $L : \iota \to$ filling and
  $Q : \iota \to \tau \to$ query algorithm, and $\varepsilon \in \mathbb{Q}$. Write
  $\mathrm{prom}_S(x) = \frac{1}{|S|} \sum_{r \in S} x(r)$ and, for a query algorithm $t$ and
  $f : \alpha \to \mathrm{Bool}$, $\mathrm{errQ}(t, f)$ for $0$ if $t$ answers $\mathrm{true}$
  exactly when $f$ has an option equal to $\mathrm{true}$ and $1$ otherwise. Assume that for every
  $f$, $\mathrm{prom}_S(r \mapsto \mathrm{errQ}(\mathrm{encadenar}(L_r, Q_r), f)) \le \varepsilon$.
  Then
  \[
    (1 - 2\varepsilon) |\alpha| \le \mathrm{prom}_S\bigl(r \mapsto \mathrm{costoLl}(L_r, \mathrm{nula}) + \mathrm{costo}(Q_r(\mathrm{tablaDe}(L_r, \mathrm{nula})), \mathrm{nula})\bigr).
  \]
  \own{}
  \leanfile{Sagitta/Map/ConservacionDual.lean}
\end{theorem}

\begin{proof}
  \leanok
  Apply the inequality of Yao's principle, in its easy direction, to the distribution that puts weight one half on $\mathrm{nula}$
  and one half on the uniform choice of an indicator function; a deterministic algorithm that asks
  $q$ options on $\mathrm{nula}$ errs with probability at least $\frac12 (1 - q/|\alpha|)$, and
  the bound is averaged over the seeds.
\end{proof}

\begin{theorem}[Some seed does at most the average cost, for weighted finite averages]
  \label{thm:information-yao}
  \lean{Sagitta.Map.Tiempo.ConservacionDual.yao_abstracto}
  \leanok
  Let $\iota$ and $\kappa$ be types with $\kappa$ finite, $S$ a nonempty finite subset of $\iota$,
  $w : \kappa \to \mathbb{Q}$ with $w(k) \ge 0$ for every $k$ and $\sum_k w(k) = 1$,
  $c : \iota \to \kappa \to \mathbb{Q}$ and $M \in \mathbb{Q}$. Assume that for every $k$,
  $\frac{1}{|S|} \sum_{r \in S} c(r, k) \le M$. Then there is $r \in S$ with
  $\sum_k w(k) c(r, k) \le M$.
  \own{}
  \leanfile{Sagitta/Map/ConservacionDual.lean}
\end{theorem}

\begin{proof}
  \leanok
  Exchange the two averages: the average over the seeds of the weighted sum is at most $M$, and
  some seed is at most the average.
\end{proof}

\section{Short and fast programs, and time as a channel}
\label{sec:information-time-channel}

\begin{definition}[Short and fast programs for a set of instances, on cslib machines]
  \label{def:information-short-fast}
  \lean{Sagitta.Map.Informacion.Tiempo.ProgramasCortos}
  \leanok
  \uses{def:information-ktmt-log}
  Let $\mathrm{enc} : \alpha \hookrightarrow \{0,1\}^*$, $V : \alpha \times \{0,1\}^* \to
  \mathrm{Bool}$ and $A \subseteq \alpha$. $\mathrm{ProgramasCortos}(A)$ says: there is $c$ such that
  every $a \in A$ is solved, in the sense of Definition~\ref{def:information-ktmt-log}, by a machine
  numbered $i$ with clock $t$, where $n = |\mathrm{enc}(a)|$, $i \le c \lfloor \log_2 (n + 2)
  \rfloor + c$ and $t \le (n + 2)^c$.
  \own{}
  \leanfile{Sagitta/Map/Informacion/Tiempo.lean}
\end{definition}

\begin{theorem}[Short and fast programs exist if and only if machine $\mathrm{kt}$ is logarithmic, on cslib machines]
  \label{thm:information-short-iff-kt}
  \lean{Sagitta.Map.Informacion.Tiempo.cortos_iff_kt}
  \leanok
  \uses{def:information-short-fast, def:information-ktmt-log}
  Let $\mathrm{enc} : \alpha \hookrightarrow \{0,1\}^*$, $V : \alpha \times \{0,1\}^* \to
  \mathrm{Bool}$ and $A \subseteq \alpha$. Then $\mathrm{ProgramasCortos}(A)$ holds if and only if
  $\mathrm{KtMTLog}(A)$ holds.
  \own{}
  \leanfile{Sagitta/Map/Informacion/Tiempo.lean}
\end{theorem}

\begin{proof}
  \leanok
  From short and fast programs, $\mathrm{ktMT}(a) \le i + \lceil \log_2 t \rceil \le 4c \lfloor
  \log_2 (n+2) \rfloor$. From a logarithmic $\mathrm{kt}$, a machine that attains it has index at
  most $c \lfloor \log_2 (n+2) \rfloor$ and a clock below $2^{c \lfloor \log_2 (n+2) \rfloor}$,
  which is at most $(n+2)^c$.
\end{proof}

\begin{theorem}[Capacity of the time channel, in a counting model]
  \label{thm:information-time-channel}
  \lean{Sagitta.Map.Informacion.Tiempo.capacidad_canal_tiempo}
  \leanok
  Let $n, T \in \mathbb{N}$ and $f : \mathrm{Fin}(2^n) \to \mathbb{N}$ be injective with $f(m) \le T$
  for every $m$. Then $2^n \le T + 1$.
  \own{}
  \leanfile{Sagitta/Map/Informacion/Tiempo.lean}
\end{theorem}

\begin{proof}
  \leanok
  Corestrict $f$ to $\mathrm{Fin}(T + 1)$ and compare cardinalities of an injection.
\end{proof}

\begin{theorem}[Search in order sends a word as its running time, in a counting model]
  \label{thm:information-search-encodes}
  \lean{Sagitta.Map.Informacion.Tiempo.busqueda_codifica}
  \leanok
  \uses{thm:information-time-channel}
  For $n \in \mathbb{N}$ and a verifier $V : \mathrm{Fin}(2^n) \to \mathrm{Bool}$, let
  $\mathrm{tiempoBusqueda}(V)$ be the number of candidates tried when the candidates $0, 1, \dots$
  are tried in order, extended by false outside the range, until the first one that $V$ accepts
  (with fuel $2^n$). Then the map $w \mapsto \mathrm{tiempoBusqueda}(x \mapsto [x = w])$ on
  $\mathrm{Fin}(2^n)$ is injective, and for every $w$ its value is $w + 1$.
  \own{}
  \leanfile{Sagitta/Map/Informacion/Tiempo.lean}
\end{theorem}

\begin{proof}
  \leanok
  The search stops exactly after the first accepted candidate, which for the verifier of equality to
  $w$ is $w$ itself, so the number of candidates tried is $w + 1$; injectivity follows.
\end{proof}

\section{Determinant and permanent modulo two}
\label{sec:information-det-perm}

\begin{theorem}[The determinant equals the permanent over the integers modulo two]
  \label{thm:information-det-perm}
  \lean{Sagitta.Map.Informacion.Tiempo.det_eq_permanent_zmod2}
  \leanok
  Let $m$ be a finite type with decidable equality and $M$ a matrix over $\mathbb{Z}/2\mathbb{Z}$
  indexed by $m \times m$. Then $\det M = \mathrm{perm}\, M$.
  \source{\cite{Valiant79}.}
  \fidelity{Not checked against the printed source.}
  \leanfile{Sagitta/Map/Informacion/Tiempo.lean}
\end{theorem}

\begin{proof}
  \leanok
  Expand both as sums over permutations; the sign of a permutation is $1$ or $-1$ and $-1 = 1$ in
  characteristic two, so the terms agree.
\end{proof}

\section{Boxes and pipes}
\label{sec:information-box}

The model is the black box of Section~\ref{sec:information-work}, with the leaves of a tree of
choices. Let $\iota$ and $\eta$ be finite types. An oracle $f : \iota \times \eta \to \mathrm{Bool}$
has the subtree of $q \in \iota$ as its leaves $(q, y)$, and $\mathrm{hayDebajo}(f, q)$ says
that some leaf $(q, y)$ of the subtree is true. A box is a map $\mathrm{caja}$ from oracles to
$\mathrm{Fin}(s) \to \mathrm{Bool}$, with no cost and no knowledge of the query that will be asked.
A run for the query $q$ is a query algorithm $R(q)$ over $\mathrm{Fin}(s) \oplus (\iota \times
\eta)$ that reads a bit of the box or a leaf at cost one.

\begin{definition}[A run that answers from a box, and that reads at most $t$; own runs]
  \label{def:information-box-run}
  \lean{Sagitta.Map.Tiempo.PlanoCaja.Responde, Sagitta.Map.Tiempo.PlanoCaja.Lee, Sagitta.Map.Tiempo.CanoProducto.Propia}
  \leanok
  \uses{def:information-filling}
  For $\mathrm{caja}$ and $R$ as above, $\mathrm{Responde}(\mathrm{caja}, R)$ says: for every $f$ and
  every $q$, the answer of $R(q)$ on the combined oracle that reads the box
  $\mathrm{caja}(f)$ on the left and $f$ on the right is $\mathrm{hayDebajo}(f, q)$. For
  $t \in \mathbb{N}$, $\mathrm{Lee}(R, t)$ says that for every $q$ and every combined oracle $g$, the
  cost of $R(q)$ on $g$ is at most $t$. $\mathrm{Propia}(R)$ says that for every $q$, every $g$ and
  every leaf $p$ read by $R(q)$ on $g$, $p$ belongs to the subtree $q$, that is, $p_1 = q$.
  \own{}
  \leanfile{Sagitta/Map/PlanoCaja.lean}
\end{definition}

\begin{theorem}[Lower bound for a static box, in a black-box model]
  \label{thm:information-static-box}
  \lean{Sagitta.Map.Tiempo.PlanoCaja.cota_caja_estatica}
  \leanok
  \uses{def:information-box-run}
  Let $\iota, \eta$ be finite types, $s, t \in \mathbb{N}$, $\mathrm{caja} : (\iota \times \eta \to
  \mathrm{Bool}) \to \mathrm{Fin}(s) \to \mathrm{Bool}$ and $R : \iota \to \text{query
  algorithm over } \mathrm{Fin}(s) \oplus (\iota \times \eta)$. Assume $\mathrm{Responde}(\mathrm{caja},
  R)$ and $\mathrm{Lee}(R, t)$. Then $(|\iota| - s) \cdot |\eta| \le 2^s \cdot |\iota| \cdot t$,
  with the subtraction of natural numbers.
  \own{}
  \leanfile{Sagitta/Map/PlanoCaja.lean}
\end{theorem}

\begin{proof}
  \leanok
  Let $Z$ be the set of leaves read by some run on the zero oracle with some content of the box;
  it has at most $2^s |\iota| t$ elements. A subtree with all its leaves in $Z$ is saturated, and in
  every other subtree there is a leaf outside $Z$; putting the ones only on such leaves, no run sees
  a one, so the box alone has to say which of the unsaturated subtrees are full, which takes
  $2^{\text{unsaturated}} \le 2^s$ contents.
\end{proof}

\begin{theorem}[Product form of the bound for a box, with own runs, in a black-box model]
  \label{thm:information-product-form}
  \lean{Sagitta.Map.Tiempo.CanoProducto.forma_producto}
  \leanok
  \uses{def:information-box-run}
  Under the hypotheses of the previous theorem, assume also $\mathrm{Propia}(R)$, $t \le |\eta|$ and
  $0 < |\eta|$. Then, in the real numbers, $|\iota| \cdot (1 - t / |\eta|) \le s$.
  \own{}
  \leanfile{Sagitta/Map/CanoProducto.lean}
\end{theorem}

\begin{proof}
  \leanok
  For own runs the count of the previous argument gives $(2|\eta| - t)^{|\iota|} \le 2^s
  |\eta|^{|\iota|}$ in the natural numbers; take logarithms and use $\log (1 + x) \ge x \log 2$ for
  $x \in [0, 1]$.
\end{proof}

\section{Description gives no lower bound}
\label{sec:information-descriptive}

\begin{theorem}[Charging the description at each step does not leave polynomial time, in a cost model]
  \label{thm:information-description-poly}
  \lean{Sagitta.Map.Tiempo.CotaDescriptiva.descripcion_no_saca_de_pol}
  \leanok
  % uses-lean: Sagitta.Map.Tiempo.Pol
  Let $L \in \mathbb{N}$ and $t, S : \mathbb{N} \to \mathbb{N}$ with $t \in \mathrm{Pol}$ and
  $S(n) \le t(n)$ for every $n$, where $\mathrm{Pol}$ is the set of functions $f$ with
  $f(n) \le c (n+1)^k$ for all $n$ for some $c, k$. Let $\mathrm{costoReal}(L, t, S)(n) = t(n)
  (L + S(n))$, a declared cost of $t(n)$ steps that each read a program of $L$ bits and a state of
  $S(n)$ bits. Then $\mathrm{costoReal}(L, t, S) \in \mathrm{Pol}$.
  \own{}
  \leanfile{Sagitta/Map/CotaDescriptiva.lean}
\end{theorem}

\begin{proof}
  \leanok
  The cost is at most $t(n) (L + t(n))$, a product and a sum of members of $\mathrm{Pol}$, and
  $\mathrm{Pol}$ is closed downward.
\end{proof}

\section{Blum speed-up}
\label{sec:information-blum}

The cost in this section is $\mathrm{cost}(c, n)$, the least \texttt{evaln} budget with which the
code $c$ halts on $n$ (a measure of Blum, not machine time). A code $c$ \emph{computes} $f$ when
$c$ on every $n$ halts with value $f(n)$. For functions $g, h : \mathbb{N} \to \mathbb{N}$,
$g = O(h)$ means that there are $C, N$ with $g(n) \le C h(n)$ for all $n \ge N$.

\begin{definition}[Speed-up for a rate $r$, on codes of Mathlib]
  \label{def:information-speedup}
  \lean{Sagitta.Map.Tiempo.AceleracionBlum.Computa, Sagitta.Map.Tiempo.AceleracionBlum.TieneAceleracion}
  \leanok
  For $r : \mathbb{N} \to \mathbb{N} \to \mathbb{N}$ and $f : \mathbb{N} \to \mathbb{N}$,
  $\mathrm{TieneAceleracion}(r, f)$ says: some code computes $f$, and for every code $i$ that
  computes $f$ there is a code $j$ that computes $f$ and a bound $N$ with
  $r(n, \mathrm{cost}(j, n)) \le \mathrm{cost}(i, n)$ for all $n \ge N$.
  \source{\cite{Blum67}.}
  \fidelity{Not checked against the printed source.}
  \leanfile{Sagitta/Map/AceleracionBlum.lean}
\end{definition}

\begin{theorem}[A function with speed-up has no optimal program, on codes of Mathlib]
  \label{thm:information-blum-sin-optimo}
  \lean{Sagitta.Map.Tiempo.AceleracionBlum.sin_optimo}
  \leanok
  \uses{def:information-speedup}
  Let $r : \mathbb{N} \to \mathbb{N} \to \mathbb{N}$ and $f : \mathbb{N} \to \mathbb{N}$ with
  $(n + 1) x \le r(n, x)$ for all $n, x$, and assume $\mathrm{TieneAceleracion}(r, f)$. Then there
  is no code $i$ that computes $f$ such that for every code $j$ that computes $f$,
  $\mathrm{cost}(i, \cdot) = O(\mathrm{cost}(j, \cdot))$.
  \source{\cite{Blum67}.}
  \fidelity{Not checked against the printed source.}
  \leanfile{Sagitta/Map/AceleracionBlum.lean}
\end{theorem}

\begin{proof}
  \leanok
  From a program $i$, speed-up with the rate bounded below by $(n+1)x$ gives a program $j$ whose
  cost is negligible against that of $i$; together with $\mathrm{cost}(j, \cdot) = O(\mathrm{cost}(i,
  \cdot))$ negligibility forces $\mathrm{cost}(j, n) = 0$ from some point on, but the cost of a
  halting code exceeds its input.
\end{proof}

\begin{theorem}[Existence of a function with speed-up for every computable rate, on codes of Mathlib]
  \label{thm:information-blum-exists}
  \lean{Sagitta.Map.Tiempo.BlumExiste.blum}
  \leanok
  \uses{def:information-speedup}
  Let $r : \mathbb{N} \to \mathbb{N} \to \mathbb{N}$ be computable as a function of two variables.
  Then there is $f : \mathbb{N} \to \mathbb{N}$ with $f(n) \le 1$ for all $n$, computable, and
  $\mathrm{TieneAceleracion}(r, f)$.
  \source{\cite{Blum67}.}
  \fidelity{Not checked against the printed source.}
  \leanfile{Sagitta/Map/BlumExiste.lean}
\end{theorem}

\begin{proof}
  \leanok
  A diagonalization with cancellation, built with the recursion theorem of Mathlib and the
  s-m-n theorem; the cost of each level is measured by an unbounded search instead of being bounded
  from the syntax. The case of a monotone rate gives the general one through its monotone envelope.
\end{proof}

\section{The cut of costs}
\label{sec:information-cut}

\begin{theorem}[A function whose programs form a strict ladder with unbounded codes, on codes of Mathlib]
  \label{thm:information-irrational-cut}
  \lean{Sagitta.Map.Tiempo.CorteRacional.existe_irracional}
  \leanok
  \uses{thm:information-blum-exists}
  There is a computable $f : \mathbb{N} \to \mathbb{N}$ with $f(n) \le 1$ for all $n$ and a sequence
  of codes $s : \mathbb{N} \to \text{codes}$ such that every $s_k$ computes $f$,
  $\mathrm{cost}(s_{k+1}, \cdot) \in o(\mathrm{cost}(s_k, \cdot))$ for every $k$, and for every $m
  \in \mathbb{N}$ there is $k$ with $m < \mathrm{encode}(s_k)$. Here $g \in o(h)$ means that for
  every $c$ there is $N$ with $c \cdot g(n) \le h(n)$ for all $n \ge N$, and $\mathrm{encode}$ is the
  natural number that Mathlib assigns to a code.
  \own{}
  \leanfile{Sagitta/Map/CorteRacional.lean}
\end{theorem}

\begin{proof}
  \leanok
  Take the function of Blum speed-up for the linear rate. The ladder descends from any program; its
  costs are positive and strictly decreasing in the order of $o$, so the programs are pairwise
  different, and there are finitely many codes below any bound.
\end{proof}

\begin{theorem}[Levin's search is below every rung of a ladder of solvers, with a free verifier]
  \label{thm:information-levin-below}
  \lean{Sagitta.Map.Tiempo.CorteRacional.levin_en_O}
  \leanok
  \uses{thm:information-closed-law}
  Let $V : \mathbb{N} \to \mathbb{N} \to \mathrm{Bool}$, $i : \mathbb{N} \to \mathbb{N}$ and
  $t : \mathbb{N} \to \mathbb{N} \to \mathbb{N}$ with $1 \le t(k, x)$ for all $k, x$, and assume
  that for all $k, x$ there is $y$ with $\mathrm{correr}(i(k), t(k, x), x) = \mathrm{some}\ y$ and
  $V(x, y) = \mathrm{true}$. Then for every $k$, $x \mapsto \mathrm{costoLevin}(V, x)$ is
  $O(x \mapsto t(k, x))$.
  \own{}
  \leanfile{Sagitta/Map/CorteRacional.lean}
\end{theorem}

\begin{proof}
  \leanok
  For each $k$ and $x$, Levin's bound for one program gives $\mathrm{costoLevin}(V, x) < 8 \cdot
  2^{i(k)} t(k, x)$, so the constant is $8 \cdot 2^{i(k)}$.
\end{proof}

\section{Writing an index}
\label{sec:information-index}

\begin{theorem}[Writing one of $m$ cases takes at least $\lceil \log_k m \rceil$ symbols]
  \label{thm:information-write-index}
  \lean{Sagitta.Map.Tiempo.EscalaLog.codificar_log}
  \leanok
  Let $m, k, \ell \in \mathbb{N}$ and $e : \mathrm{Fin}(m) \to (\mathrm{Fin}(\ell) \to
  \mathrm{Fin}(k))$ be injective. Then $\lceil \log_k m \rceil \le \ell$.
  \own{}
  \infrastructure{}
  \leanfile{Sagitta/Map/EscalaLog.lean}
\end{theorem}

\begin{proof}
  \leanok
  An injection between finite types gives $m \le k^\ell$ by cardinality, and the ceiling of the
  logarithm is the least exponent with that property.
\end{proof}
